\section{Split And Tree Certificates}

A split instance binds its tree path, parent-certificate identity, canonical
unit universe, populations, weighted edges, seat counts, and orientation rule.
The result is either an optimal connected assignment or exact infeasibility.
The bounded reference oracle commits to every candidate in deterministic order.

A whole-plan certificate must establish more than the validity of individual
cuts. The verifier reconstructs both child contexts from each certified parent
assignment. Submitted child populations, edges, and unit IDs are not trusted.
Non-leaf nodes must appear in the breadth-first order of the canonical
bisection schedule. One-seat leaves are ordered by binary path and must
partition the root universe exactly once.

This construction distinguishes algorithmic proof from plan audit. An RPLAN
audit can confirm assignment shape, population, contiguity, and provenance.
The recursive certificate establishes the stronger model-scoped claim that no
lexicographically better cut exists at each certified node.

The guarantee remains sequential. If a locally optimal parent cut produces a
child that cannot be split under the same rules, the procedure fails
explicitly; it does not silently backtrack to a worse parent cut.
